% Part: turing-machines
% Chapter: undecidability
% Section: unsolvability-decision-problem

\documentclass[../../../include/open-logic-section]{subfiles}

\begin{document}

\olfileid{tur}{und}{uns}
\olsection{నిర్ణయ సమస్య పరిష్కరించలేనిది}

\begin{thm}
\ollabel{thm:decision-prob}
నిర్ణయ సమస్య పరిష్కరించలేనిది: మొదటిస్థాయి తర్కంలోని \tetoken{ఒక వాక్యం}{!!a{sentence}}~$!B$ గల
టేపుపై దాన్ని ఇన్‌పుట్‌గా తీసుకుని $D$ను ప్రారంభించినప్పుడు చివరకు ఆగి, $!B$
చెల్లుబాటైతే అప్పుడే మాత్రమే $1$ను, లేకపోతే $0$ను అవుట్‌పుట్‌గా ఇచ్చే
ట్యూరింగ్ యంత్రం~$D$ లేదు.
\end{thm}

\begin{proof}
నిర్ణయ సమస్య పరిష్కరించదగినదని, అంటే అలాంటి ట్యూరింగ్ యంత్రం~$D$ ఉందని
అనుకుందాం. అప్పుడు ఆగే సమస్యను కింది విధంగా పరిష్కరించగలిగేవారం. ట్యూరింగ్
యంత్రం~$M_e$ సంఖ్య~$e$, ఇన్‌పుట్~$w$లను ఇన్‌పుట్‌గా ఇచ్చినప్పుడు, అనుగుణమైన
\tetoken{వాక్యం}{!!{sentence}}~$!T(M_e,w) \lif !E(M_e,w)$ను గణించి, టేపుపై అత్యంత ఎడమ గడిని
చదువుతూ ఆగే ట్యూరింగ్ యంత్రం~$E$ను నిర్మిస్తాం. అప్పుడు యంత్రం
$E\concat D$, ఇన్‌పుట్ $e$,~$w$లను ఇచ్చినప్పుడు మొదట
$!T(M_e,w) \lif !E(M_e,w)$ను గణించి, తరువాత ఆ ఇన్‌పుట్‌పై నిర్ణయ-సమస్య
యంత్రం~$D$ను నడుపుతుంది. $!T(M_e,w) \lif !E(M_e,w)$ చెల్లుబాటైతే
అప్పుడే మాత్రమే $D$ అవుట్‌పుట్~$1$తో ఆగుతుంది; లేకపోతే అవుట్‌పుట్~$0$ను
ఇస్తుంది. \olref[ver]{lem:halt-if-valid}, \olref[ver]{lem:valid-if-halt}
ప్రకారం, $M_e$ ఇన్‌పుట్~$w$పై ఆగితే, అప్పుడే మాత్రమే
$!T(M_e,w) \lif !E(M_e,w)$ చెల్లుబాటవుతుంది. కాబట్టి $E\concat D$కు
ఇన్‌పుట్ $e$,~$w$లను ఇచ్చినప్పుడు, $M_e$ ఇన్‌పుట్~$w$పై ఆగితే అప్పుడే
మాత్రమే అది అవుట్‌పుట్~$1$తో ఆగుతుంది; లేకపోతే అవుట్‌పుట్~$0$తో ఆగుతుంది.
మరో మాటలో, $E\concat D$ ఆగే సమస్యను పరిష్కరించేది. కానీ
\olref[hal]{thm:halting-problem} ప్రకారం అలాంటి ట్యూరింగ్ యంత్రం ఉండదని
మనకు తెలుసు.
\end{proof}

\begin{cor}\ollabel{cor:undecidable-sat}%
మొదటిస్థాయి తర్కంలోని ఏకపక్ష \tetoken{వాక్యం}{!!{sentence}} సంతృప్తిపరచదగినదో కాదో నిర్ణయించలేం.
\end{cor}

\begin{proof}
  సంతృప్తిపరచదగినతను ట్యూరింగ్ యంత్రం~$S$ నిర్ణయిస్తుందని అనుకుందాం.
  అప్పుడు నిర్ణయ సమస్యను కింది విధంగా పరిష్కరించగలిగేవారం:
  \tetoken{ఒక వాక్యం}{!!a{sentence}}~$B$ను ఇన్‌పుట్‌గా ఇచ్చినప్పుడు $!B$ను ఒక గడి కుడివైపుకు
  తరలించండి. గడి~$1$కు తిరిగి వెళ్లి $\lnot$ సంకేతాన్ని వ్రాయండి.

  ఇప్పుడు ట్యూరింగ్ యంత్రం~$S$ను నడపండి. $\lnot !B$ సంతృప్తిపరచదగినదైతే
  $1$తో, $\lnot !B$ సంతృప్తిపరచలేనిదైతే~$0$తో, అది చివరకు టేపుపై ఆగుతుంది. గడి~$1$పై
  ఒక~$\TMstroke$ ఉంటే దాన్ని చెరపండి; గడి~$1$ ఖాళీగా ఉంటే ఒక~$\TMstroke$ను
  వ్రాసి, తరువాత ఆగండి.

  ఈ ట్యూరింగ్ యంత్రం ఎల్లప్పుడూ ఆగుతుంది; $\lnot !B$ సంతృప్తిపరచలేనిదైతే
  అప్పుడే మాత్రమే దాని అవుట్‌పుట్~$1$, లేకపోతే~$0$. $!B$ చెల్లుబాటైతే,
  అప్పుడే మాత్రమే $\lnot !B$ సంతృప్తిపరచలేనిది కనుక, $!B$ చెల్లుబాటైతే
  అప్పుడే మాత్రమే యంత్రం~$1$ను, లేకపోతే~$0$ను అవుట్‌పుట్‌గా ఇస్తుంది;
  అంటే అది నిర్ణయ సమస్యను పరిష్కరించేది.
\end{proof}

\begin{explain}
అందువల్ల ``$!B$ మొదటిస్థాయి తర్కంలోని చెల్లుబాటయ్యే \tetoken{వాక్యం}{!!{sentence}}నా?'' అనే
ప్రశ్నకు ఎల్లప్పుడూ సరైన ``అవును'' లేదా ``కాదు'' సమాధానం ఇచ్చే ట్యూరింగ్
యంత్రం లేదు. అయితే, ఎల్లప్పుడూ సరైన ``అవును'' సమాధానం ఇచ్చే---కానీ సమాధానం
``కాదు'' అయితే ఆగనే ఆగని---ట్యూరింగ్ యంత్రం \emph{ఉంది}. ఇది మొదటిస్థాయి
తర్కపు నిర్దుష్టత, సంపూర్ణత సిద్ధాంతాల నుండీ, \tetoken{వ్యుత్పత్తులను}{!!{derivation}s} ప్రభావకంగా
లెక్కించవచ్చనే వాస్తవం నుండీ అనుగమిస్తుంది.
\end{explain}

\begin{thm}
  \ollabel{thm:valid-ce}%
  మొదటిస్థాయి \tetoken{వాక్యాల}{!!{sentence}s} చెల్లుబాటుతనం అర్ధ-నిర్ణయించదగినది:
  మొదటిస్థాయి తర్కంలోని \tetoken{ఒక వాక్యం}{!!a{sentence}}~$!B$ గల టేపుపై ట్యూరింగ్ యంత్రం~$E$ను
  ఇన్‌పుట్‌తో ప్రారంభించినప్పుడు, $!B$ చెల్లుబాటైతే అప్పుడే మాత్రమే చివరకు
  ఆగి అవుట్‌పుట్~$1$ను ఇచ్చే, లేకపోతే ఆగని ట్యూరింగ్ యంత్రం~$E$ ఉంది.
\end{thm}

\begin{proof}
  మొదటిస్థాయి తర్కపు సాధ్యమైన అన్ని \tetoken{వ్యుత్పత్తులను}{!!{derivation}s} ఒక ప్రభావక
  అల్గోరిథంతో ఒకదాని తరువాత మరొకటిగా ఉత్పత్తి చేయవచ్చు. యంత్రం~$E$ ఇదే
  చేస్తుంది; $\Proves !B$ అని చూపే \tetoken{ఒక వ్యుత్పత్తిని}{!!a{derivation}} కనుగొన్నప్పుడు
  అవుట్‌పుట్~$1$తో ఆగుతుంది. నిర్దుష్టత సిద్ధాంతం ప్రకారం, $E$ అవుట్‌పుట్~$1$తో
  ఆగితే, దానికి కారణం~$\Entails !B$. సంపూర్ణత సిద్ధాంతం ప్రకారం,
  $\Entails !B$ అయితే $\Proves !B$ అని చూపే \tetoken{ఒక వ్యుత్పత్తి}{!!a{derivation}} ఉంటుంది.
  సాధ్యమైన అన్ని \tetoken{వ్యుత్పత్తులను}{!!{derivation}s} $E$ క్రమబద్ధంగా ఉత్పత్తి చేస్తుంది కనుక,
  $\Proves !B$ అని చూపేదాన్ని అది చివరకు కనుగొని, అవుట్‌పుట్~$1$తో ఆగుతుంది.
\end{proof}

\end{document}
